% All thesis-specific details live here, edit this file and nothing else in
% the front matter. Every value is consumed by dimscthesis.cls, which builds
% the cover, the examination committee page and the abstracts from them.

% Author. The template prints the name as "Όνομα Μ. Επώνυμο", where Μ. is the
% initial of the father's first name, and the abbreviated form in the footer
% of every page of the main body.
\authorFirstGr{Χρήστος}
\authorFirstAbrGr{Χ.}
\authorMiddleGr{Δ.}
\authorLastGr{Μουρτογιάννης}
\authorFirstEn{Christos}
\authorFirstAbrEn{C.}
\authorMiddleEn{D.}
\authorLastEn{Mourtogiannis}

% Αριθμός μητρώου, printed on the examination committee page.
\authorSN{ME1.20.0005}

% The interdepartmental postgraduate programme this thesis is submitted to,
% run jointly by the Department of Informatics and Telecommunications and the
% Department of Economics.
\programGr{%
  ΔΙΑΤΜΗΜΑΤΙΚΟ ΠΡΟΓΡΑΜΜΑ ΜΕΤΑΠΤΥΧΙΑΚΩΝ ΣΠΟΥΔΩΝ
  «ΔΙΟΙΚΗΣΗ ΚΑΙ ΟΙΚΟΝΟΜΙΚΗ ΤΩΝ ΤΗΛΕΠΙΚΟΙΝΩΝΙΑΚΩΝ ΔΙΚΤΥΩΝ
  ΚΑΙ ΠΛΗΡΟΦΟΡΙΑΚΩΝ ΣΥΣΤΗΜΑΤΩΝ»%
}
\programEn{%
  INTERDEPARTMENTAL POSTGRADUATE PROGRAMME IN MANAGEMENT AND ECONOMICS
  OF TELECOMMUNICATION NETWORKS AND INFORMATION SYSTEMS%
}

% Titles. The Greek one is the title as registered on the department's
% declaration form, character for character.
\titleGr{Εξαγωγή Ευφυών Αναλύσεων από Δεδομένα Ηλεκτρονικού Εμπορίου}
\titleEn{Extraction of Intelligent Analytics from E-Commerce Data}

% Month and year printed on the cover. The template requires the month and
% year of the examination, not of the writing.
\dateGr{ΣΕΠΤΕΜΒΡΙΟΣ 2026}
\dateEn{SEPTEMBER 2026}

% Full examination date, printed on the examination committee page.
\examinationDateGr{30 Σεπτεμβρίου 2026}
\examinationDateEn{September 30, 2026}

% Supervisor, who is also the first member of the examination committee.
\advisorGr{Ευστάθιος Χατζηευθυμιάδης}
\advisorRankGr{Καθηγητής}
\advisorOrgGr{ΕΚΠΑ}
\advisorEn{Efstathios Hadjiefthymiades}
\advisorRankEn{Professor}
\advisorOrgEn{NKUA}

% Second member of the examination committee.
\boardTwoGr{Κωνσταντίνος Κολομβάτσος}
% The class lays the committee out as one line per member beside a wide
% label, and this member's qualification does not fit beside it. The break
% goes after the name, so the rank and the institution stay together.
\boardTwoRankGr{\\Επίκουρος Καθηγητής}
\boardTwoOrgGr{Πανεπιστήμιο Θεσσαλίας}
\boardTwoEn{Konstantinos Kolomvatsos}
\boardTwoRankEn{Assistant Professor}
\boardTwoOrgEn{University of Thessaly}

% Third member of the examination committee.
\boardThreeGr{Μιλτιάδης Κυριακάκος}
% Broken after the name for the same reason as the second member above.
\boardThreeRankGr{\\Εργαστηριακό Διδακτικό Προσωπικό}
\boardThreeOrgGr{ΕΚΠΑ}
\boardThreeEn{Miltiadis Kyriakakos}
\boardThreeRankEn{Laboratory Teaching Staff}
\boardThreeOrgEn{NKUA}

% Supervisors block on the cover and the committee page. A single supervisor
% keeps the default heading, ΕΠΙΒΛΕΠΩΝ.
\supervisorsGr{%
  {\bfseries Ευστάθιος Χατζηευθυμιάδης}, Καθηγητής ΕΚΠΑ%
}
\supervisorsEn{%
  {\bfseries Efstathios Hadjiefthymiades}, Professor NKUA%
}

% Subject area and keywords, printed after each abstract. The template allows
% up to five keywords.
\subjectAreaGr{Εξόρυξη Δεδομένων και Επιχειρηματική Ευφυΐα}
\keywordsGr{εξόρυξη δεδομένων, ηλεκτρονικό εμπόριο, συχνά υποσύνολα,
  συσταδοποίηση, επιχειρηματική ευφυΐα}
\subjectAreaEn{Data Mining and Business Intelligence}
\keywordsEn{data mining, electronic commerce, frequent itemsets, clustering,
  business intelligence}
